\documentclass[11pt,a4paper]{article}
\usepackage[margin=1in]{geometry}
\usepackage{amsmath,amssymb,amsthm}
\usepackage{physics}
\usepackage{hyperref}
\usepackage{booktabs}
\usepackage{enumitem}
\usepackage{tcolorbox}
\usepackage{bm}

\newcommand{\rhat}{\hat{r}}
\newcommand{\that}{\hat{\theta}}
\newcommand{\zhat}{\hat{z}}
\newcommand{\xhat}{\hat{x}}
\newcommand{\yhat}{\hat{y}}
\newcommand{\HT}{\mathcal{H}}
\newcommand{\E}{\mathbf{E}}
\newcommand{\J}{\mathbf{J}}

\title{Cylindrical Far-Field Propagation (CyFFP)\\[6pt]
\large Complete Formulation: Vector Cylindrical Harmonics,\\
FFTLog Hankel Transforms, Graf's Addition Theorem,\\
and the Neumann Shift Theorem}
\author{Arvin Keshvari \quad (Supervisor: Dr.\ Zin Lin)}
\date{March 2026}

\begin{document}
\maketitle
\tableofcontents
\newpage

%=======================================================================
\section{Introduction and Problem Statement}
%=======================================================================

We compute the point spread function (PSF) of rotationally symmetric optical systems
(metalenses, ideal lenses) at oblique incidence angles $\alpha$.
The near field on a cylindrical monitor plane is decomposed into
angular-momentum eigenmodes $e^{im\theta}$, propagated to the focal plane via
vector cylindrical harmonics, and re-expanded in a local basis centered at the
focal spot $(f\tan\alpha,\,0,\,f)$.

The formulation is \emph{exact} (no paraxial approximation) and handles arbitrarily
large oblique angles, provided sufficient angular modes are retained.

%=======================================================================
\section{Vector Cylindrical Harmonics}
\label{sec:harmonics}
%=======================================================================

\subsection{Near-field decomposition}

The tangential electric field on the near-field plane $z = 0$ in cylindrical
coordinates $(r,\theta)$ is
\begin{equation}\label{eq:nearfield}
  \E^{\mathrm{Near}}(r,\theta)
  = \sum_m \bigl[E_{m,r}(r)\,\rhat + E_{m,\theta}(r)\,\that\bigr]\,e^{im\theta},
\end{equation}
where the modal amplitudes are obtained by FFT over $\theta$:
\begin{equation}
  E_{m,r/\theta}(r_j)
  = \frac{1}{N_\theta}\sum_l E_{r/\theta}(r_j,\theta_l)\,e^{-im\theta_l}.
\end{equation}

\subsection{TE harmonics}

From the Helmholtz equation in cylindrical coordinates, the TE electric eigenmode is
\begin{equation}\label{eq:TE}
  \E^{\mathrm{TE}}_{m,k_r}(r,\theta,z)
  = -\Bigl[\rhat\,\frac{im}{r}\,J_m(k_r r)
         - \that\,\frac{\partial J_m(k_r r)}{\partial r}\Bigr]
    \,e^{im\theta}\,e^{ik_z z},
\end{equation}
where $k_z = \sqrt{k^2 - k_r^2}$ and $k = 2\pi/\lambda$.

\subsection{TM harmonics}

The tangential part of the TM electric eigenmode is
\begin{equation}\label{eq:TM}
  \E^{\mathrm{TM},\perp}_{m,k_r}
  = -\frac{k_z}{k}\,\frac{\partial J_m}{\partial r}\,\rhat
    \;-\; i\,\frac{m k_z}{k\, r}\,J_m(k_r r)\,\that.
\end{equation}

\subsection{Bessel recurrence identities}

The NIST recurrence relations (used throughout):
\begin{align}
  J_{m+1}(x) + J_{m-1}(x) &= \frac{2m}{x}\,J_m(x), \label{eq:rec1}\\
  J_{m-1}(x) - J_{m+1}(x) &= 2\,\frac{dJ_m(x)}{dx}, \label{eq:rec2}\\
  J_{-m}(x) &= (-1)^m\,J_m(x). \label{eq:negm}
\end{align}


%=======================================================================
\section{TE and TM Expansion Coefficients}
\label{sec:coeffs}
%=======================================================================

Projecting the near field onto the TE/TM eigenmodes and using the orthogonality
of Bessel functions $\int_0^\infty r\,J_m(k_r r)\,J_m(k_r' r)\,dr = \delta(k_r - k_r')/k_r$,
we obtain:

\begin{tcolorbox}[title=TE coefficient]
\begin{equation}\label{eq:ATE}
  A^{\mathrm{TE}}_m(k_r)
  = \frac{i}{2}\bigl[\HT_{m+1}[r\,E_{m,r}] + \HT_{m-1}[r\,E_{m,r}]\bigr]
  + \frac{1}{2}\bigl[\HT_{m-1}[r\,E_{m,\theta}] - \HT_{m+1}[r\,E_{m,\theta}]\bigr]
\end{equation}
\end{tcolorbox}

\begin{tcolorbox}[title=TM coefficient]
\begin{equation}\label{eq:ATM}
  A^{\mathrm{TM}}_m(k_r)
  = -\frac{k_z/k}{2}\bigl[\HT_{m-1}[r\,E_{m,r}] - \HT_{m+1}[r\,E_{m,r}]\bigr]
  - \frac{i\,k_z/k}{2}\bigl[\HT_{m-1}[r\,E_{m,\theta}] + \HT_{m+1}[r\,E_{m,\theta}]\bigr]
\end{equation}
\end{tcolorbox}

\noindent
where $\HT_\nu[f](k_r) = \int_0^\infty f(r)\,J_\nu(k_r r)\,dr$ is the order-$\nu$
Hankel transform with the plain $dr$ measure.  Each mode~$m$ requires
\textbf{4~FFTLog calls} (2 if $E_{m,\theta} = 0$).

The mode truncation is $|m| \le M_{\max} = \lceil k\sin\alpha\cdot R\rceil + \delta M$,
determined by the Bessel decay $J_m(x) \approx 0$ for $m \gg x$.


%=======================================================================
\section{Negative-$m$ Symmetry}
\label{sec:symmetry}
%=======================================================================

For $x$-polarized illumination ($\varphi = 0$) of an axisymmetric structure, the
source current satisfies $\J_{-m} = \hat\sigma\,\J_m$ with
$\hat\sigma = \mathrm{diag}(1,-1,1)$.  Maxwell's equations then give
\begin{equation}\label{eq:sym}
  E_{-m,r} = E_{m,r}, \qquad E_{-m,\theta} = -E_{m,\theta}.
\end{equation}
Substituting into~\eqref{eq:ATE}--\eqref{eq:ATM} and using
$\HT_{-\nu}[f] = (-1)^\nu\,\HT_\nu[f]$ (from $J_{-\nu}(x)=(-1)^\nu J_\nu(x)$):

\begin{tcolorbox}[title=Symmetry relations]
\begin{equation}\label{eq:symTE}
  A^{\mathrm{TE}}_{-m} = (-1)^{m+1}\,A^{\mathrm{TE}}_m, \qquad
  A^{\mathrm{TM}}_{-m} = (-1)^{m}\,A^{\mathrm{TM}}_m.
\end{equation}
\end{tcolorbox}

\noindent
Combined propagated coefficient:
\begin{equation}\label{eq:Atilde_neg}
  \tilde A_{-m}(k_r) = (-1)^m\bigl[A^{\mathrm{TM}}_m(k_r) - A^{\mathrm{TE}}_m(k_r)\bigr]\cdot e^{ik_z f}.
\end{equation}

\textbf{Consequence:} Only $m = 0,1,\ldots,M_{\max}$ need be computed; negative modes
are reconstructed algebraically.  FDTD/RCWA simulations need only run for $m \ge 0$.


%=======================================================================
\section{FFTLog Hankel Transform}
\label{sec:fftlog}
%=======================================================================

FFTLog (Talman 1978; Hamilton 2000) computes $\HT_\nu[f](k_r)$ for \emph{any}
real order~$\nu$ by a log-variable change $r = r_0 e^{-\rho}$,
$k_r = k_0 e^{\kappa}$, converting the Hankel integral into a convolution
evaluated by FFT.  The filter kernel in Fourier space is
\begin{equation}\label{eq:fftlog_kernel}
  \hat{\tilde J}_\nu(q)
  = \frac{\Gamma\!\bigl(\tfrac{\nu+1+q}{2}\bigr)}
         {\Gamma\!\bigl(\tfrac{\nu+1-q}{2}\bigr)}\,2^q,
\end{equation}
computed via \texttt{loggamma} for overflow-free evaluation at $|\nu| \sim 600$.

\textbf{There is no turning-point problem:} FFTLog is order-agnostic, unlike
asymptotic methods (e.g.\ Beckman--O'Neil NUFHT) which break down at
$k_r r \approx |m|$.

\paragraph{Grid constraint.}
The input must be on a log-spaced grid $r_j = r_0\,e^{j\,\Delta\ln}$.
The output lives on the reciprocal log-grid with
$k_r^{\max} = 1/r_{\min}$.  Therefore
\begin{equation}
  r_{\min} \;\le\; \frac{1}{k} = \frac{\lambda}{2\pi}
\end{equation}
is \textbf{required} to capture all propagating modes.  For a 2\,mm lens at
$\lambda = 500$\,nm this gives $r_{\min} \le 80$\,nm and
$N_r \sim 2^{17} = 131072$.


%=======================================================================
\section{Propagation to the Focal Plane}
\label{sec:propagation}
%=======================================================================

Combine TE and TM, apply the propagation phase:
\begin{equation}\label{eq:prop}
  \tilde A_m(k_r)
  = \bigl[A^{\mathrm{TE}}_m(k_r) + A^{\mathrm{TM}}_m(k_r)\bigr]\,e^{ik_z(k_r)\,f},
  \quad k_z = \sqrt{k^2 - k_r^2}.
\end{equation}
Evanescent modes ($k_r > k$, imaginary $k_z$) are set to zero.


%=======================================================================
\section{Graf's Addition Theorem: Local-Basis Shift}
\label{sec:graf}
%=======================================================================

\subsection{The theorem}

For a shift $\bm{x}_0 = (x_0, 0)$ with $x_0 = f\tan\alpha$, define local polar
coordinates $(\rho,\psi)$ by $x = x_0 + \rho\cos\psi$, $y = \rho\sin\psi$.
Graf's addition theorem states:
\begin{equation}\label{eq:graf}
  J_n(k_r r)\,e^{in\theta}
  = \sum_{m=-\infty}^{\infty} J_m(k_r x_0)\,J_{n-m}(k_r\rho)\,e^{i(n-m)\psi},
  \qquad \rho < x_0.
\end{equation}

\subsection{Local modal coefficients}

Substituting~\eqref{eq:graf} into the far-field sum and relabelling $l = n - m$:
\begin{tcolorbox}[title=Graf shift (mode convolution)]
\begin{equation}\label{eq:Bl}
  B_l(k_r)
  = \sum_{|m|\le M_{\max}} \tilde A_{m+l}(k_r)\,J_m(k_r\,x_0),
  \qquad |l| \le L_{\max}.
\end{equation}
\end{tcolorbox}

\noindent
This is a discrete convolution in mode index at each $k_r$.
The weight $J_m(k_r x_0)$ decays exponentially for $|m| > k_r x_0$,
so the effective sum length is $\min(M_{\max},\,\lceil k_r x_0\rceil + 20)$.

\paragraph{Normal incidence.}
For $\alpha \to 0$: $x_0 \to 0$, $J_m(0) = \delta_{m,0}$, so $B_l = \tilde A_l$
(identity).

\paragraph{Comparison with polar-to-Cartesian regridding.}
The earlier approach (cyFFP1) assembled $\mathcal{U}(k_r,\phi)$ via IFFT over $m$
and interpolated onto a Cartesian $(\tilde k_x, \tilde k_y)$ grid.
The Graf approach is \emph{spectrally exact}: Eq.~\eqref{eq:Bl} is an algebraic
identity with exponentially small truncation error for $\rho_{\max} < x_0$.

\subsection{Implementation: Miller backward recurrence}

Computing $J_0(x),\ldots,J_{M_{\max}}(x)$ via Miller backward recurrence is
$\sim$10$\times$ faster than independent \texttt{besselj} calls.
Negative orders use $J_{-m}(x) = (-1)^m J_m(x)$.  The inner accumulation loop
uses \texttt{@simd} for SIMD vectorization.


%=======================================================================
\section{Inverse Hankel Transforms in the Local Basis}
\label{sec:inverse}
%=======================================================================

The field in the local frame is a sum of cylindrical harmonics:
\begin{equation}\label{eq:bl}
  b_l(\rho)
  = \int_0^\infty B_l(k_r)\,J_l(k_r\rho)\,k_r\,dk_r
  = \HT_l\bigl[k_r\,B_l(k_r)\bigr](\rho).
\end{equation}
Since the Hankel transform is self-inverse (with $k_r\,dk_r$ measure), this is
another FFTLog call of order~$l$.  Only $2L_{\max}+1$ calls are needed, with
\begin{equation}
  L_{\max} \approx k \cdot \mathrm{NA} \cdot \rho_{\max}.
\end{equation}
For $\mathrm{NA} = 0.5$, $\rho_{\max} = 5\lambda$: $L_{\max} \approx 16$.
Compare $M_{\max} \approx 600$--$4000$ for the forward transforms.


%=======================================================================
\section{Angular Synthesis}
\label{sec:synthesis}
%=======================================================================

The PSF in local polar coordinates is
\begin{equation}\label{eq:psf}
  E(\rho_j, \psi_s, f)
  = \sum_{l=-L_{\max}}^{L_{\max}} b_l(\rho_j)\,e^{il\psi_s},
\end{equation}
evaluated by IFFT over $l$ at each radial point $\rho_j$.
Take $|E|^2$ for the PSF intensity.


%=======================================================================
\section{Neumann Shift Theorem: Scalar Fast Path}
\label{sec:neumann}
%=======================================================================

For near fields of the LPA form
$\E(r,\theta) = t(r)\,e^{ik_x r\cos\theta}\,(\cos\theta\,\rhat - \sin\theta\,\that)$,
we can bypass the $M_{\max}$ individual FFTLog calls entirely.

\subsection{Scalar modal coefficients}

The Jacobi--Anger expansion gives
$u_m(r) = 2\pi(-i)^m\,t(r)\,J_m(k_x r)$,
and the scalar spectral coefficient is
\begin{equation}
  a_m(k_r) = k_r \int_0^\infty t(r)\,J_m(k_x r)\,J_m(k_r r)\,r\,dr.
\end{equation}

\subsection{Neumann addition formula}

The product of same-order Bessel functions satisfies
\begin{equation}\label{eq:neumann}
  J_m(a\,r)\,J_m(b\,r)
  = \frac{1}{2\pi}\int_0^{2\pi}
    J_0\!\bigl(c(\varphi)\,r\bigr)\,e^{-im\varphi}\,d\varphi,
  \qquad c(\varphi) = \sqrt{a^2 + b^2 - 2ab\cos\varphi}.
\end{equation}

Substituting into $a_m$:
\begin{tcolorbox}[title=Neumann shift theorem]
\begin{equation}\label{eq:neumann_shift}
  a_m(k_r) = (-i)^m\,k_r \cdot
  \frac{1}{2\pi}\int_0^{2\pi}
    \underbrace{T_0\!\bigl(c(k_r, k_x, \varphi)\bigr)}_{\text{single Hankel transform}}
    \;e^{-im\varphi}\,d\varphi,
\end{equation}
\end{tcolorbox}
\noindent
where
\begin{equation}
  T_0(k) = \HT_0[r\,t(r)](k) = \int_0^\infty r\,t(r)\,J_0(k\,r)\,dr
\end{equation}
is a \textbf{single order-0 FFTLog call}, and $c(\varphi) = \sqrt{k_r^2 + k_x^2 - 2k_r k_x\cos\varphi}$.

\paragraph{Algorithm.}
\begin{enumerate}[nosep]
  \item Compute $T_0(k)$ via one FFTLog call.
  \item For each $k_r$ (threaded): evaluate $T_0(c(\varphi))$ at $N_\varphi$ angles
        by interpolation on the log-spaced $k_r$ grid.
  \item FFT over $\varphi$ to extract $a_m$ for all $|m| \le M_{\max}$ simultaneously.
\end{enumerate}

\paragraph{Cost.}
$O(N_r \log N_r)$ for the single FFTLog, plus $O(N_{k_r} \cdot N_\varphi \log N_\varphi)$
for the interpolation and FFTs.  This replaces $O(M_{\max} \cdot N_r \log N_r)$ FFTLog
calls in the standard path --- a reduction of $\sim M_{\max}/1 \sim 4000\times$ in the
Hankel transform step.


%=======================================================================
\section{Complete Algorithm}
\label{sec:algorithm}
%=======================================================================

\begin{tcolorbox}[title=CyFFP: Standard Path]
\textbf{Inputs:} $E_r$, $E_\theta$ on log-spaced $r$, uniform $\theta$;\;
$k$, $\alpha$, $f$.

\medskip
\textbf{Step 1} --- Angular FFT (over $\theta$): extract $E_{m,r/\theta}$
for $m = 0,\ldots,M_{\max}$.

\textbf{Step 2} --- Forward Hankel transforms (FFTLog, orders $m\pm 1$):
$A^{\mathrm{TE}}_m$, $A^{\mathrm{TM}}_m$ for $m \ge 0$ only.

\textbf{Step 3} --- Propagation + symmetry:
$\tilde A_m = (A^{\mathrm{TE}}_m + A^{\mathrm{TM}}_m)\,e^{ik_z f}$;\;
$\tilde A_{-m} = (-1)^m(A^{\mathrm{TM}}_m - A^{\mathrm{TE}}_m)\,e^{ik_z f}$.

\textbf{Step 4} --- Graf shift:
$B_l(k_r) = \sum_m \tilde A_{m+l}(k_r)\,J_m(k_r x_0)$,\;
$|l| \le L_{\max}$.

\textbf{Step 5} --- Inverse Hankel (FFTLog, order $l$):
$b_l(\rho) = \HT_l[k_r B_l](\rho)$.

\textbf{Step 6} --- Angular IFFT (over $l$):
$E(\rho,\psi) = \sum_l b_l(\rho)\,e^{il\psi}$.
\end{tcolorbox}

\begin{tcolorbox}[title=CyFFP: Neumann Shift Fast Path (scalar/LPA fields)]
\textbf{Input:} radial transmission $t(r)$;\; $k$, $\alpha$, $f$.

\medskip
\textbf{Step A} --- Single FFTLog call: $T_0(k) = \HT_0[r\,t(r)](k)$.

\textbf{Step B} --- For each $k_r$ (threaded): evaluate
$T_0(\sqrt{k_r^2 + k_x^2 - 2k_r k_x\cos\varphi})$, multiply by
$(-i)^m k_r\,e^{ik_z f}$, FFT over $\varphi \to \tilde A_m(k_r)$ for all $m$.

\textbf{Steps 4--6} --- Unchanged (Graf shift, inverse Hankel, synthesis).
\end{tcolorbox}


%=======================================================================
\section{Oblique Ideal Lens and the Airy Disk}
\label{sec:airy}
%=======================================================================

The ideal oblique lens phase that focuses perfectly at $(f\tan\alpha,\,0,\,f)$ is
\begin{equation}\label{eq:oblique_lens}
  u(r,\theta) = \exp\!\Bigl(-ik\Bigl[
    \sqrt{(r\cos\theta - f\tan\alpha)^2 + r^2\sin^2\!\theta + f^2}
    \;-\; f\Bigr]\Bigr).
\end{equation}
This produces a perfect converging spherical wave with no aberrations at the
design angle.  The PSF is an Airy disk with effective NA
\begin{equation}\label{eq:NA_eff}
  \mathrm{NA}_{\mathrm{eff}}
  = \frac{R\cos\alpha}{\sqrt{R^2\cos^2\!\alpha + f^2}},
\end{equation}
which reduces to $R/\sqrt{R^2+f^2}$ at normal incidence.  The Airy first zero is at
$\rho_1 = 3.8317/(k\cdot\mathrm{NA}_{\mathrm{eff}})$.

Note that $\mathrm{NA}_{\mathrm{eff}} < \mathrm{NA}_0$: the off-axis focal spot is at
distance $f/\cos\alpha > f$ from the aperture centre, so the aperture subtends a
smaller half-angle and the Airy disk is wider than at normal incidence.

This should \emph{not} be confused with a normal-incidence lens illuminated by an
oblique plane wave, $t(r)\cdot e^{ik\sin\alpha\,r\cos\theta}$, which has coma
aberrations and a non-ideal PSF.


%=======================================================================
\section{Performance Optimizations}
\label{sec:perf}
%=======================================================================

\begin{center}
\begin{tabular}{@{}lll@{}}
\toprule
\textbf{Optimization} & \textbf{Speedup} & \textbf{Mechanism} \\
\midrule
Negative-$m$ symmetry & $2\times$ on Step~2 & Only $m \ge 0$ computed \\
$E_\theta = 0$ detection & $2\times$ on Step~2 & 2 FFTLog calls instead of 4 \\
Neumann shift (LPA) & ${\sim}M_{\max}\times$ on Step~2 & 1 FFTLog instead of $M_{\max}$ \\
Bessel truncation & ${\sim}4\times$ on Step~4 & $m_{\mathrm{cut}} = \min(M_{\max},\,\lceil k_r x_0\rceil+20)$ \\
Miller recurrence & ${\sim}10\times$ on Bessel eval & $J_0\ldots J_{M_{\max}}$ in one $O(M_{\max})$ pass \\
\texttt{@simd} inner loop & $2$--$4\times$ on Step~4 & SIMD vectorization of dot product \\
\texttt{Threads.@threads} & ${\sim}N_{\mathrm{threads}}\times$ & Steps 2, 4, 5, 6 parallelized \\
\bottomrule
\end{tabular}
\end{center}


%=======================================================================
\section{Comparison of Approaches}
\label{sec:comparison}
%=======================================================================

\begin{center}
\begin{tabular}{@{}lp{5.5cm}p{5.5cm}@{}}
\toprule
& \textbf{cyFFP1} (polar$\to$Cartesian) & \textbf{cyFFP2} (Graf shift) \\
\midrule
Steps 3--6 & IFFT over $m$; interpolate; 2D IFFT
           & Propagation; mode convolution; FFTLog$_l$; IFFT over $l$ \\
Interpolation & Bilinear (accuracy-limited) & None --- spectrally exact \\
Output grid & Cartesian $(\xi,\eta)$ & Polar $(\rho,\psi)$ local \\
Modes needed & $\sim(k\,\mathrm{NA}\,\rho_{\max})^2$ Cartesian pts
             & $L_{\max} \sim k\,\mathrm{NA}\,\rho_{\max}$ \\
Accuracy vs $\alpha$ & Degrades (larger grid) & Uniform (exponential convergence) \\
\bottomrule
\end{tabular}
\end{center}


%=======================================================================
\section*{References}
%=======================================================================

\begin{enumerate}[label={[\arabic*]},nosep]
\item J.~D.~Jackson, \emph{Classical Electrodynamics}, 3rd ed.\ (Wiley, 1999).
\item M.~Abramowitz and I.~A.~Stegun, \emph{Handbook of Mathematical Functions}
      (Dover, 1965).
\item F.~W.~J.~Olver et al., \emph{NIST Digital Library of Mathematical Functions},
      \url{https://dlmf.nist.gov/}.
\item J.~D.~Talman, ``Numerical Fourier and Bessel transforms in logarithmic variables,''
      \emph{J.\ Comput.\ Phys.}\ \textbf{29}, 35--48 (1978).
\item A.~J.~S.~Hamilton, ``Uncorrelated modes of the non-linear power spectrum,''
      \emph{Mon.\ Not.\ R.\ Astron.\ Soc.}\ \textbf{312}, 257--284 (2000).
\item DLMF \S10.23(ii): Graf's addition theorem for Bessel functions.
\item H.~Goldstein, C.~Poole, and J.~Safko, \emph{Classical Mechanics}, 3rd ed.\
      (Addison-Wesley, 2002).
\end{enumerate}

\end{document}
